To provide a rigorous grounding for our proposed framework, we situate our work at the intersection of Multi-Agent Reinforcement Learning (MARL), communicative coordination, and the Free Energy Principle (FEP).

\subsection{Multi-Agent Reinforcement Learning and Non-Stationarity}
Standard MARL frameworks typically rely on the Centralized Training and Decentralized Execution (CTDE) paradigm to mitigate the non-stationarity arising from concurrent agent updates \cite{lowe2017multi, rashid2018qmix}. While value-decomposition methods like QMIX \cite{rashid2018qmix} and actor-critic variants like MAPPO \cite{yu2022surprising} have shown success in static or cooperative benchmarks, they often struggle in highly dynamic, competitive-cooperative domains such as Real-Time Bidding (RTB) \cite{zhao2024adaptive}. In these environments, the Markov assumption is frequently violated as agents' policies shift in response to market volatility \cite{papoudakis2021benchmarking}. Recent attempts to address non-stationarity have focused on meta-learning \cite{al-shedivat2018continuous} and opponent modeling \cite{xie2024learning}. However, these approaches often converge to sub-optimal equilibria or facilitate implicit collusion—where agents inadvertently coordinate to inflate prices or suppress competition—due to a purely instrumental focus on reward maximization \cite{calvano2020artificial, schut2024collusion}. Our work diverges by shifting the objective from reward-seeking to variational surprise minimization, which provides a principled regularization against such collusive artifacts.

\subsection{Communication in Multi-Agent Coordination}
The emergence of communication in MAS has been traditionally framed as a reinforcement problem where agents learn to transmit signals that maximize a shared return \cite{foerster2016learning, sukhbaatar2016learning}. Methods such as DIAL \cite{foerster2016learning} and TarMAC \cite{das2019tarmac} utilize differentiable channels to propagate gradients across agents. More recent work has explored sparse communication to reduce bandwidth overhead \cite{wang2024survey} and the use of Large Language Models (LLMs) as zero-shot coordinators \cite{shinn2024reflexion, zhang2024proagent}. Despite these advances, most existing protocols treat communication as an instrumental action—a means to an external reward \cite{hu2024towards}. This often leads to ``cheap talk'' or the failure of communication protocols when reward structures are perturbed \cite{sharma2017learning}. In contrast, our framework treats communication as an \textit{epistemic action} \cite{parrill_epistemic_2023}. By optimizing for the reduction of uncertainty about the latent states of other agents and the environment, our agents communicate to build shared ``world models,'' a concept rooted in the theory of minds and collective active inference \cite{friston2024collective}.

\subsection{Active Inference and the Free Energy Principle}
Active Inference (ActInf) offers a unified normative theory for perception and action, positing that biological agents act to minimize their variational free energy (VFE) \cite{friston2010free}. In single-agent settings, ActInf has been shown to naturally balance exploration (epistemic value) and exploitation (instrumental value) \cite{friston2017active, safron2024active}. Extensions to multi-agent settings have begun to explore ``social'' active inference, where agents model each other as intentional entities \cite{albarracin2023epistemic, smith2024active}. However, scaling these models to high-dimensional, non-stationary environments like RTB remains a challenge. Current literature has applied ActInf to robotic swarm coordination \cite{stulp2024active} and economic games \cite{friedman2024active}, but these often assume pre-defined state spaces. Our contribution bridges this gap by integrating hierarchical generative models with LLM-augmented inference \cite{wei2022chain}, allowing agents to communicate and reason over novel, unencoded concepts in real-time. This aligns with recent shifts toward ``Foundational MARL,'' where agents leverage large-scale pre-training to handle environmental complexity \cite{yang2024foundation}.

\subsection{The Frontier: Epistemic vs. Instrumental Agency}
The critical distinction of our approach lies in the prioritization of \textit{epistemic} over \textit{instrumental} agency. While modern MARL baselines like QPLEX \cite{wang2021qplex} or FACMAC \cite{peng2021facmac} optimize for the ``What'' (the reward), ActInf optimizes for the ``Why'' (the underlying causal structure). Recent 2024 studies in AI Safety highlight that purely reward-driven agents are prone to goal misgeneralization \cite{langosco2024goal}. By grounding communication in the minimization of Expected Free Energy (EFE), our framework ensures that information exchange is intrinsically motivated by the need for accurate world-modeling, providing a robust defense against the non-stationarity and coordination failures prevalent in state-of-the-art systems \cite{mousavi2024comprehensive, duan2024multiagent}.